\section{Graph}
\subsection{Graph Condense (Kosaraju SCC)}
\begin{minted}{c++}
struct SCC {
  int n, comps;
  vector<vector<int>> adj, rev_adj, scc_adj;
  vector<int> order, comp;
  vector<bool> vis;
  SCC(int _n) : n(_n), comps(0), adj(_n), rev_adj(_n),
    comp(_n, -1), vis(_n, false) {}
  void add_edge(int u, int v) {
    adj[u].push_back(v);
    rev_adj[v].push_back(u);
  }
  void dfs1(int u) {
    vis[u] = true;
    for (int v : adj[u]) if (!vis[v]) dfs1(v);
    order.push_back(u);
  }
  void dfs2(int u, int id) {
    comp[u] = id;
    for (int v : rev_adj[u]) if (comp[v] == -1) dfs2(v, id);
  }
  void build() {
    for (int i = 0; i < n; i++) if (!vis[i]) dfs1(i);
    reverse(order.begin(), order.end());
    for (int u : order) {
      if (comp[u] == -1) dfs2(u, comps++);
    }
    scc_adj.resize(comps);
    for (int u = 0; u < n; u++) {
      for (int v : adj[u]) {
        if (comp[u] != comp[v]) {
          scc_adj[comp[u]].push_back(comp[v]);
        }
      }
    }
  }
};
\end{minted}
\subsection{Articulation Points \texorpdfstring{\&}{\&} Bridges}
\begin{minted}{c++}
struct Articulation {
  int n, timer;
  vector<vector<int>> adj;
  vector<int> tin, low;
  vector<bool> is_art;
  vector<pair<int, int>> bridges;
  Articulation(int _n) : n(_n), timer(0), adj(_n),
    tin(_n, -1), low(_n, -1), is_art(_n, false) {}
  void add_edge(int u, int v) {
    adj[u].push_back(v);
    adj[v].push_back(u);
  }
  void dfs(int u, int p = -1) {
    tin[u] = low[u] = timer++;
    int children = 0;
    for (int v : adj[u]) {
      if (v == p) continue;
      if (tin[v] != -1) {
        low[u] = min(low[u], tin[v]);
      } else {
        children++;
        dfs(v, u);
        low[u] = min(low[u], low[v]);
        if (low[v] >= tin[u] && p != -1) is_art[u] = true;
        if (low[v] > tin[u]) bridges.push_back({min(u,v), max(u,v)});
      }
    }
    if (p == -1 && children > 1) is_art[u] = true;
  }
  void build() {
    for (int i = 0; i < n; i++) {
      if (tin[i] == -1) dfs(i);
    }
  }
};
\end{minted}
\subsection{Hopcroft-Karp (Bipartite Matching)}
\begin{minted}{c++}
struct HopcroftKarp { // Bipartite Matching
  int n, m;
  vector<vector<int>> g;
  vector<int> dist, pairU, pairV;
  static const int INF = 1e9;
  HopcroftKarp(int n_ = 0, int m_ = 0) {
    n = n_; m = m_;
    g.assign(n, {});
    pairU.assign(n, -1);
    pairV.assign(m, -1);
    dist.assign(n, 0);
  }
  void add_edge(int u, int v) {
    g[u].push_back(v);
  }
  bool bfs() {
    queue<int> q;
    for (int u = 0; u < n; ++u) {
      if (pairU[u] == -1) {
        dist[u] = 0;
        q.push(u);
      } else {
        dist[u] = INF;
      }
    }
    bool foundAug = false;
    while (!q.empty()) {
      int u = q.front(); q.pop();
      for (int v : g[u]) {
        int u2 = pairV[v];
        if (u2 == -1) {
          foundAug = true;
        } else if (dist[u2] == INF) {
          dist[u2] = dist[u] + 1;
          q.push(u2);
        }
      }
    }
    return foundAug;
  }
  bool dfs(int u) {
    for (int v : g[u]) {
      int u2 = pairV[v];
      if (u2 == -1 || (dist[u2] == dist[u] + 1 &&
      dfs(u2))) {
        pairU[u] = v;
        pairV[v] = u;
        return true;
      }
    }
    dist[u] = INF;
    return false;
  }
  int max_matching() {
    int matching = 0;
    while (bfs()) {
      for (int u = 0; u < n; ++u)
        if (pairU[u] == -1 && dfs(u))
          ++matching;
    }
    return matching;
  }
  vector<pair<int,int>> matching_edges() const
  {
    vector<pair<int,int>> res;
    for (int v = 0; v < m; ++v)
      if (pairV[v] != -1)
        res.emplace_back(pairV[v], v);
    return res;
  }
};
\end{minted}
\subsection{Dinic's Algorithm (Max Flow)}
\begin{minted}{c++}
const int64_t inf = 1LL << 61;
struct Dinic {
  struct edge {
    int to, rev;
    int64_t flow, w;
  };
  int n, s, t;
  vector<int> d;
  vector<int> done;
  vector<vector<edge>> g;
  Dinic() {};
  Dinic(int _n) {
    n = _n + 10, g.resize(n);
  }
  void add_edge(int u, int v, int64_t w) {
    edge a = {v, (int)g[v].size(), 0, w};
    edge b = {u, (int)g[u].size(), 0, 0};
    g[u].push_back(a), g[v].push_back(b);
  }
  bool bfs() {
    d.assign(n, -1), d[s] = 0;
    queue<int> Q;
    Q.push(s);
    while (Q.size()) {
      int u = Q.front(); Q.pop();
      for (auto &e : g[u]) {
        int v = e.to;
        if (!~d[v] && e.flow < e.w) d[v] = d[u] + 1,
        Q.push(v);
      }
    }
    return d[t] != -1;
  }
  int64_t dfs(int u, int64_t flow) {
    if (u == t) return flow;
    for (int &i = done[u]; i < (int)g[u].size(); i++) {
      edge &e = g[u][i];
      if (e.w <= e.flow) continue;
      int v = e.to;
      if (d[v] == d[u] + 1) {
        int64_t nw = dfs(v, min(flow, e.w - e.flow));
        if (nw > 0) {
          e.flow += nw;
          g[v][e.rev].flow -= nw;
          return nw;
        }
      }
    }
    return 0;
  }
  int64_t max_flow(int _s, int _t) {
    s = _s, t = _t;
    int64_t flow = 0;
    while (bfs()) {
      done.assign(n, 0);
      while (int64_t nw = dfs(s, inf)) flow += nw;
    }
    return flow;
  }
};
\end{minted}
\subsection{Biconnected Components \texorpdfstring{\&}{\&} Bridge Tree}
\begin{minted}{c++}
struct BCC {
  int n, timer, comps;
  vector<vector<int>> adj, tree;
  vector<int> tin, low, comp;
  vector<bool> is_bridge, vis;
  vector<pair<int, int>> edges;
  BCC(int _n) : n(_n), timer(0), comps(0), adj(_n), tin(_n, -1),
    low(_n, -1), comp(_n, -1), vis(_n, false) {}
  void add_edge(int u, int v) {
    adj[u].push_back(edges.size());
    adj[v].push_back(edges.size());
    edges.push_back({u, v});
    is_bridge.push_back(false);
  }
  void dfs(int u, int p = -1) {
    vis[u] = true;
    tin[u] = low[u] = timer++;
    for (int id : adj[u]) {
      int v = edges[id].first ^ edges[id].second ^ u;
      if (v == p) continue;
      if (vis[v]) low[u] = min(low[u], tin[v]);
      else {
        dfs(v, u);
        low[u] = min(low[u], low[v]);
        if (low[v] > tin[u]) is_bridge[id] = true;
      }
    }
  }
  void build_tree() {
    for (int i = 0; i < n; i++) if (!vis[i]) dfs(i);
    vector<int> q;
    for (int i = 0; i < n; i++) {
      if (comp[i] != -1) continue;
      q.push_back(i); comp[i] = comps++;
      while (!q.empty()) {
        int u = q.back(); q.pop_back();
        for (int id : adj[u]) {
          int v = edges[id].first ^ edges[id].second ^ u;
          if (comp[v] == -1 && !is_bridge[id]) {
            comp[v] = comp[u]; q.push_back(v);
          }
        }
      }
    }
    tree.resize(comps);
    for (int i = 0; i < edges.size(); i++) {
      if (is_bridge[i]) {
        int u = comp[edges[i].first], v = comp[edges[i].second];
        tree[u].push_back(v); tree[v].push_back(u);
      }
    }
  }
};
\end{minted}
